We document state variation in Medicaid coverage for obesity-indicated GLP-1 medications over time, and use a stacked difference-in-differences design to estimate the effects of coverage on utilization and net-of-rebate spending. Nine quarters out, coverage increases prescriptions for obesity-indicated GLP-1 medications by \rxObSevenNineEmw{} per 100 enrollee-months (SE = \rxObSevenNineEmwSE{}). Coverage had no effect on GLP-1 prescribing for diabetes or cardiovascular indications, suggesting that off-label prescribing of diabetes formulations for obesity is not very common in the Medicaid program. The expansions do not appear to affect consumer spending at major online GLP-1 compounding firms, which suggests that the utilization response in our main analysis reflects new utilization rather than crowd-out. We find that coverage increases net-of-rebate Medicaid spending by \$\netSpObSevenNineEmw{} per 100 enrollee-months (SE = \$\netSpObSevenNineEmwSE{}), compared to \$\spObSevenNineEmw{} in gross reimbursements---manufacturer rebates reduce the state's net spending by approximately \spReductionObSevenNineEmw{} percent. Across the 17 states that ever covered obesity-indicated GLP-1 medications by mid 2025, coverage increased Medicaid net spending by around \$\netCostCoveringEtNine{} per year. Extending Medicaid coverage in the remaining states would generate a further \$\netCostNoncoveringEtNine{} in annual net spending.